% Vietnamese translation for OI Public Library
% Source-PDF: IOI中国国家候选队论文/2007/day2/2.古楠《平面嵌入》.pdf
\documentclass[11pt]{article}
\usepackage{oipl-vn}

\newcommand{\Oh}{\mathcal{O}}
\newcommand{\GP}{G_P}
\newcommand{\DFI}{\mathrm{DFI}}
\newcommand{\Kfive}{K_5}
\newcommand{\Kthree}{K_{3,3}}

\title{Nhúng phẳng}
\author{Gu Nan\\Trường Trung học Nanshan, Miên Dương, Tứ Xuyên}
\date{Trại đông tin học 2007}

\begin{document}
\maketitle

\origtitle{Nhúng phẳng}
\sourcepath{IOI China candidate team papers / 2007 / day 2 / Gu Nan}

\begin{abstract}
Bài viết gồm hai phần chính. Phần thứ nhất giới thiệu thuật toán nhúng phẳng:
các thao tác cụ thể, phân tích độ phức tạp, chứng minh tính đúng đắn và một số
bài liên quan. Phần thứ hai là phụ lục, gồm tài liệu tham khảo và giả mã.
\end{abstract}

\paragraph{Từ khóa.}
Mặt phẳng; nhúng; hoạt động trong; liên quan; hoạt động ngoài; khối; cạnh gốc;
cạnh hồi.

\section{Thuật toán}

\subsection{Các định nghĩa liên quan}

Trước hết nhắc lại vài khái niệm cơ bản.

\begin{enumerate}
  \item \term{Vẽ phẳng}: một đồ thị có thể được biến đổi thành hình vẽ mà mọi
  cạnh không cắt nhau, trừ giao nhau tại đỉnh. Quá trình biến đổi ấy gọi là vẽ
  phẳng. Những đỉnh vốn kề nhau vẫn phải kề nhau sau khi biến đổi.

  \item \term{Đồ thị phẳng}: đồ thị có thể vẽ phẳng được gọi là đồ thị phẳng;
  ngược lại là đồ thị không phẳng.

  \item \term{Nhúng phẳng}: tương đương với vẽ phẳng, nhưng được mô tả dưới
  dạng lưu trữ: với mỗi đỉnh, ta lưu danh sách các đỉnh kề theo thứ tự chiều
  kim đồng hồ. Bảng dùng để lưu gọi là danh sách kề.
\end{enumerate}

Trong bài này, mọi lần nói đến ``mặt phẳng'' đều hiểu theo các định nghĩa trên,
không phải mặt phẳng hình học tùy ý.

\begin{center}
  \small Hình 1: ví dụ một đồ thị được vẽ lại sao cho các cạnh không cắt nhau.
\end{center}

\subsection{Mục tiêu của thuật toán}

Với một đồ thị \(G\) có \(n\) đỉnh và \(m\) cạnh, mục tiêu là dùng thời gian
\(\Oh(n)\) để phán đoán đồ thị có phẳng hay không. Nếu đồ thị phẳng, thuật toán
cũng phải dựng được một nhúng phẳng trong thời gian \(\Oh(n)\).

\subsection{Kiến thức chuẩn bị}

Ta cần các kiến thức như tìm kiếm theo chiều sâu, thành phần song liên thông,
đỉnh khớp và sắp xếp đếm. Ngoài ra cần một vài định lý cơ bản: số cạnh của
một đồ thị phẳng không vượt quá \(3n-5\); nếu vượt ngưỡng đó thì đồ thị chắc
chắn không phẳng. Kuratowski chứng minh rằng hai đồ thị cản trở nhúng phẳng là
\(\Kfive\) và \(\Kthree\). Một đồ thị chứa một đồ thị đồng phôi với \(\Kfive\)
hoặc \(\Kthree\) thì không phẳng, và chiều ngược lại cũng đúng.

Thao tác cơ bản của thuật toán là thêm cạnh. Ta tạo một đồ thị \(\GP\), rồi
đưa các cạnh của đồ thị gốc vào \(\GP\) theo một thứ tự nhất định. Thao tác đó
được gọi là \term{nhúng cạnh}.

\begin{center}
  \small Hình 2: hai chướng ngại cổ điển cho tính phẳng, \(\Kfive\) và
  \(\Kthree\).
\end{center}

\subsection{Tổng quan thuật toán}

Để trình bày rõ ràng, trước hết chỉ xét nhúng phẳng trên đồ thị song liên
thông. Với đồ thị có đỉnh khớp, ta có thể tách đồ thị tại các đỉnh khớp, nhúng
từng phần rồi ghép lại.

Trước tiên thực hiện một lần DFS, thu được cây DFS và thứ tự DFS đảo ngược
\(\DFI\). Cạnh nối một đỉnh với con của nó gọi là \term{cạnh cây}; cạnh nối
một đỉnh với hậu duệ không phải con trực tiếp gọi là \term{cạnh hồi}.

Sau đó xử lý các đỉnh theo thứ tự \(\DFI\) đảo ngược. Khi xử lý một đỉnh, ta
trước hết nhúng các cạnh cây của nó vào \(\GP\), rồi nhúng mọi cạnh hồi của
nó. Các cạnh nối với tổ tiên chưa cần xử lý. Nếu quá trình nhúng thất bại, ta
kết luận \(G\) không phẳng; nếu không, ta hoàn thành nhúng phẳng.

\subsection{Nhúng cạnh}

Nhúng cạnh là thao tác quan trọng nhất. Giả sử cần nhúng cạnh \((v,w)\). Ban
đầu, thành phần liên thông đang xét chứa một đỉnh khớp \(r\). Nếu bỏ \(r\),
thành phần này tách thành hai phần. Ta gắn một bản sao của \(r\) cho mỗi phần,
nhúng cạnh \((v,w)\), rồi hợp nhất hai thành phần lại; lúc đó \(r\) không còn
là đỉnh khớp trong khối mới.

\begin{center}
  \small Hình 3: tách một khối tại đỉnh khớp, nhúng cạnh mới, có thể đảo một
  phía rồi hợp nhất lại.
\end{center}

Trong quá trình nhúng, hễ \(\GP\) chứa đỉnh khớp thì ta tách thành phần liên
thông của nó thành hai phần. Khi nhúng cạnh hồi, các phần này lại được ghép.
Vì vậy \(\GP\) thực chất được tạo bởi nhiều thành phần song liên thông. Khi cài
đặt, đây không chỉ là cách hiểu mà là cấu trúc cần được xử lý thật sự.

\subsection{Hoạt động ngoài, liên quan và hoạt động trong}

Phép đảo được dùng để bảo vệ một tính chất quan trọng: trong quá trình nhúng
cạnh, mọi đỉnh hoạt động ngoài phải nằm trên mặt ngoài.

Mặt ngoài là mặt tiếp xúc với miền vô hạn bên ngoài đồ thị. Với một thành phần
song liên thông chỉ có một cạnh, ta xem nó như một chu trình suy biến để thao
tác trên mặt ngoài vẫn thống nhất.

\begin{center}
  \small Hình 4: cách xem một cạnh đơn như một chu trình ngoài suy biến.
\end{center}

Định nghĩa đỉnh hoạt động ngoài phức tạp hơn. Trong một thành phần song liên
thông, đỉnh có thứ tự \(\DFI\) nhỏ nhất được gọi là \term{đỉnh gốc} của thành
phần. Do phép tách chỉ sinh ra khi nhúng cạnh cây, khi một đỉnh khớp bị tách
thành nhiều bản sao, bản sao nằm cùng phần với một con sẽ trở thành gốc của
phần đó. Những đỉnh này gọi là \term{điểm sao chép}; bản duy nhất không phải
gốc gọi là \term{điểm gốc ban đầu}.

Giả sử \(r\) là gốc của một thành phần song liên thông \(B\). Khi đó trong
\(B\), \(r\) chỉ có một con. Nếu \(r\) có hai con \(c_1,c_2\), do thứ tự DFS,
khi duyệt từ \(c_1\) ta có thể đi trong \(B\) tới \(c_2\) mà không qua \(r\),
mâu thuẫn với việc \(c_2\) là con của \(r\). Gọi con duy nhất ấy là \(c\);
cạnh \(r c\) gọi là \term{cạnh gốc}, còn thành phần tương ứng gọi là
\term{khối}. Khối gắn với bản sao gốc được gọi là \term{khối con} của điểm gốc
ban đầu.

Khi đang xử lý đỉnh \(v\), một đỉnh đã xử lý được gọi là \term{hoạt động
ngoài} nếu thỏa ít nhất một trong hai điều kiện:
\begin{enumerate}
  \item Nó có cạnh trực tiếp tới một tổ tiên của \(v\).
  \item Trong một khối con của nó tồn tại đỉnh hoạt động ngoài.
\end{enumerate}

\begin{center}
  \small Hình 5: các đỉnh hình vuông là hoạt động ngoài; tính hoạt động được
  truyền ngược qua các khối con.
\end{center}

Để lấy nhanh thông tin này, với mỗi đỉnh ta lưu:
\[
  \mathit{ecpoint}(x)=
  \text{độ sâu tổ tiên sớm nhất mà \(x\) nối trực tiếp tới},
\]
\[
  \mathit{lowpoint}(x)=
  \text{độ sâu tổ tiên sớm nhất mà \(x\) hoặc hậu duệ của \(x\) nối tới}.
\]
Mỗi đỉnh còn có một danh sách hai chiều \(\mathit{SDlist}\), lưu các giá trị
\(\mathit{lowpoint}\) của các con theo thứ tự tăng dần. Vì giá trị này bị chặn
bởi \(n\), có thể dùng sắp xếp đếm để xây dựng mọi danh sách trong thời gian
tuyến tính. Khi cần biết một đỉnh có hoạt động ngoài hay không, chỉ cần kiểm
tra \(\mathit{ecpoint}\) và phần tử đầu của \(\mathit{SDlist}\) có nhỏ hơn độ
sâu của đỉnh đang xử lý hay không. Khi một con được hợp nhất vào cha, ta xóa
giá trị tương ứng khỏi \(\mathit{SDlist}\) của cha.

Tương tự, một đỉnh \term{liên quan} khi đang xử lý \(v\) nếu:
\begin{enumerate}
  \item Nó có cạnh trực tiếp tới \(v\).
  \item Trong một khối con của nó tồn tại đỉnh liên quan.
\end{enumerate}
Một đỉnh có thể vừa hoạt động ngoài vừa liên quan. Đỉnh \term{hoạt động trong}
là đỉnh liên quan nhưng không hoạt động ngoài; đỉnh \term{không hoạt động} là
đỉnh không thuộc hai loại trên. Tương ứng có khối hoạt động ngoài, khối liên
quan, khối hoạt động trong và khối không hoạt động. Hàm \texttt{walkup} sẽ thu
thập các thông tin này.

\subsection{Phép đảo}

Mục đích của phép đảo là giữ mọi đỉnh hoạt động ngoài trên mặt ngoài trong
suốt quá trình nhúng. Trong ví dụ ở hình 3, vì một đỉnh như \(y\) là hoạt động
ngoài, phần dưới phải được đảo trước khi nhúng cạnh.

Cách trực tiếp là đảo toàn bộ danh sách kề của mọi đỉnh trong khối, nhưng như
vậy quá chậm. Ta dùng cạnh gốc: mọi cạnh cây ban đầu mang giá trị \(+1\). Khi
một khối cần đảo, ta đổi giá trị cạnh gốc của nó thành \(-1\). Một đỉnh có bị
đảo hay không phụ thuộc vào số cạnh \(-1\) trên đường cây từ gốc tới nó; số đó
lẻ thì cần đảo, chẵn thì không. Sau khi xử lý xong mọi đỉnh, một DFS cuối cùng
sẽ khôi phục toàn bộ thông tin đảo.

Vì các đỉnh được xử lý theo thứ tự \(\DFI\) đảo ngược và mọi đỉnh hoạt động
ngoài nằm trên mặt ngoài, mỗi lần thêm cạnh hồi đều diễn ra trên mặt ngoài. Do
đó khi xử lý một đỉnh, ta chỉ cần duyệt mặt ngoài của các khối. Mỗi đỉnh trên
mặt ngoài lưu hai đỉnh kề ngoài của nó. Khi một khối bị đảo, ta chỉ cần hoán
đổi thông tin kề ngoài thích hợp; không cần biết tuyệt đối đang đi theo chiều
kim đồng hồ hay ngược chiều, chỉ cần biết đi vào từ phía nào thì đi ra phía
còn lại.

\begin{center}
  \small Hình 6: phép đảo đổi hướng một khối để đỉnh hoạt động ngoài vẫn nằm
  trên mặt ngoài sau khi ghép.
\end{center}

\subsection{Phân tích toàn bộ quá trình}

Khởi tạo \(\GP\) là đồ thị rỗng. Sau đó xử lý từng đỉnh theo thứ tự \(\DFI\)
đảo ngược. Khi xử lý \(v\), với mỗi cạnh hồi \((v,w)\), ta đi từ \(w\) dọc mặt
ngoài lên tới \(v\) để ghi lại thông tin liên quan; thao tác này là
\texttt{walkup}. Sau đó đi từ \(v\) xuống để nhúng cạnh; thao tác này là
\texttt{walkdown}. Nếu cuối cùng có cạnh hồi chưa được nhúng, đồ thị không
phẳng; ngược lại thu được nhúng phẳng.

Khi cần nhúng cạnh \((v,w)\), do danh sách kề lưu theo chiều kim đồng hồ, nếu
đang duyệt theo chiều kim đồng hồ thì thêm \(w\) vào đầu danh sách kề của
\(v\); nếu ngược chiều thì thêm vào cuối. Các cạnh cây ban đầu được thêm thống
nhất vào đầu danh sách.

\begin{verbatim}
DFS(G), tinh lowpoint cho moi dinh
Khoi tao G_P, SDlist va sap xep cac SDlist
for v theo thu tu DFI dao nguoc:
    for moi con c cua v:
        nhung canh goc (v, c) vao G_P
    for moi canh hoi (v, w):
        Walkup(G_P, v, w)
    for moi con c cua v:
        Walkdown(G_P, c)
    for moi canh hoi (v, w):
        if (v, w) chua nam trong G_P:
            return (khong phang, G_P)
RecoverPlanarEmbedding(G_P)
return (phang, G_P)
\end{verbatim}

\subsection{\texttt{walkup}}

\texttt{walkup} dùng để thu thập thông tin khi xử lý đỉnh \(v\). Với cạnh hồi
\((v,w)\), ta đi từ \(w\) dọc mặt ngoài lên tổ tiên \(v\). Mỗi đỉnh có một
danh sách \(\mathit{proots}\) ghi các khối con liên quan. Khi đi từ một gốc
bản sao lên điểm gốc ban đầu \(r\), ta thêm bản sao đó vào
\(\mathit{proots}(r)\); khi đó \(r\) trở thành đỉnh liên quan.

Từ \(w\) có thể đi theo hai hướng trên mặt ngoài. Hai đường này có thể chênh
lệch rất lớn, nên thuật toán duyệt đồng thời hai hướng; nhiều nhất sau hai lần
độ dài đường ngắn hơn sẽ tới gốc khối. Ngoài ra, các cạnh hồi khác nhau có thể
tạo ra phần đường trùng nhau. Vì thế mỗi đỉnh có một nhãn \(\mathit{visited}\).
Khi xử lý \(v\), mọi đỉnh đã thăm được gán \(\mathit{visited}=v\); nếu gặp lại
đỉnh như vậy thì dừng phần duyệt hiện tại.

\begin{center}
  \small Hình 7: \texttt{walkup} đi đồng thời theo hai hướng và dừng khi gặp
  gốc khối hoặc phần đã thăm.
\end{center}

\subsection{\texttt{walkdown}}

\texttt{walkdown} dùng thông tin liên quan để nhúng các cạnh hồi. Nó cũng thao
tác trên mặt ngoài, nhưng khác với \texttt{walkup}: khi duyệt xuống, thông
thường không được đi qua đỉnh hoạt động ngoài. Ngoại lệ duy nhất là khi đỉnh
đó vẫn còn là đỉnh liên quan và ta cần đi xuống khối con của nó.

Khi xử lý \(v\), ta đi xuống từ từng khối con theo hai hướng: thuận và ngược.
Nếu gặp một đỉnh có \(\mathit{proots}\) không rỗng thì đi xuống khối con liên
quan của nó, nhưng phải chọn lại hướng.

Nếu trong quá trình đi xuống gặp một đỉnh hoạt động ngoài nhưng không liên
quan, ta dừng; đỉnh đó gọi là \term{đỉnh dừng}. Để tránh dừng quá sớm, thứ tự
duyệt phải tuân theo hai quy tắc:

\begin{enumerate}
  \item Khi có nhiều khối con cần đi xuống, luôn đi xuống khối hoạt động trong
  trước; sau khi quay lại mới xử lý các khối còn lại. Vì khối hoạt động trong
  có \(\mathit{lowpoint}=v\), còn khối hoạt động ngoài có
  \(\mathit{lowpoint}<v\), chỉ cần đưa các khối có \(\mathit{lowpoint}=v\) vào
  đầu danh sách \(\mathit{proots}\).

  \item Khi đi xuống một khối liên quan và hoạt động ngoài, trước hết chọn
  hướng trực tiếp tới đỉnh hoạt động trong; nếu không có thì chọn hướng trực
  tiếp tới đỉnh vừa liên quan vừa hoạt động ngoài. Khi gặp đỉnh dừng hoặc gốc
  khối thì dừng.
\end{enumerate}

Khi đi tới một đỉnh \(w\), các trường hợp được xét theo thứ tự:
\begin{enumerate}
  \item Nếu \(w\) là đỉnh liên quan, trước hết nhúng cạnh \((v,w)\).
  \item Nếu \(w\) không hoạt động ngoài, đi tiếp tới đỉnh kế tiếp.
  \item Nếu \(w\) hoạt động ngoài nhưng không liên quan, dừng lần duyệt này.
  \item Nếu \(w\) vừa liên quan vừa là đỉnh khớp, đi xuống khối con của nó.
\end{enumerate}

\begin{center}
  \small Hình 8: sau khi nhúng cạnh tới \(w\), nếu \(w\) vẫn còn liên quan thì
  tiếp tục đi xuống khối con; nếu không, \(w\) trở thành đỉnh dừng.
\end{center}

Để hợp nhất các thành phần song liên thông sau khi nhúng cạnh, khi đi xuống
một khối con ta đẩy lên ngăn xếp bốn thông tin: điểm không sao chép, điểm sao
chép, hướng đi xuống và hướng quay lại. Khi hợp nhất, nếu hai hướng mâu thuẫn
thì cần đảo khối trước khi ghép.

Còn một tối ưu hóa quan trọng. Nếu gặp đỉnh dừng \(s\), đoạn mặt ngoài từ một
đỉnh hoạt động trong \(p\) tới \(s\) thường không còn ích lợi nhưng vẫn có thể
bị duyệt lại. Ta thêm một cạnh giả \((v,s)\) để bỏ qua đoạn này trong các lần
duyệt tiếp theo. Mỗi khối song liên thông nhiều nhất gặp hai đỉnh dừng, nên
tổng số cạnh giả không vượt quá \(2n\); cuối thuật toán chỉ cần xóa chúng.

\begin{center}
  \small Hình 9: cạnh giả nối tới đỉnh dừng giúp tránh duyệt lại một đoạn mặt
  ngoài không còn cần thiết.
\end{center}

\section{Độ phức tạp}

DFS và tính \(\mathit{lowpoint}\) hiển nhiên là tuyến tính. Việc xây dựng mọi
\(\mathit{SDlist}\) cũng tốn \(\Oh(n)\), và mỗi lần xóa phần tử khỏi danh sách
này là hằng số nên tổng vẫn là \(\Oh(n)\). Phép đảo chỉ xảy ra khi hợp nhất
thành phần song liên thông, nhiều nhất \(n\) lần, mỗi lần hằng số. Nhúng mỗi
cạnh cũng là hằng số.

Trong \texttt{walkup} và \texttt{walkdown}, các đoạn mặt ngoài đã đi qua sẽ bị
cập nhật bởi cạnh được nhúng và không bị duyệt lại. Khi duyệt khối hoạt động
trong tuy không thêm cạnh giả, khối ấy cũng không bị duyệt lần thứ hai vì nó
không có cạnh nối với tổ tiên sớm hơn. Do đó mỗi cạnh chỉ bị duyệt số lần hằng
số. Vì đồ thị phẳng có không quá \(3n-5\) cạnh, cộng thêm \(2n\) cạnh giả vẫn
là tuyến tính. Tổng độ phức tạp của toàn bộ thuật toán là \(\Oh(n)\).

\section{Chứng minh đúng đắn}

Nếu thuật toán sai thì mọi cải tiến về độ phức tạp đều vô nghĩa, nên phần này
rất quan trọng.

Ngoài ngoại lệ đã nêu, \texttt{walkdown} không đi qua đỉnh hoạt động ngoài; vì
vậy khi nhúng cạnh nó không che mất các đỉnh hoạt động ngoài. Nếu mọi đỉnh
hoạt động ngoài luôn nằm trên mặt ngoài, các cạnh mới không cắt nhau.

Với đồ thị không phẳng, ta chỉ cần chỉ ra trong mỗi tình huống thất bại đều có
một đồ thị con đồng phôi với \(\Kfive\) hoặc \(\Kthree\).

\paragraph{Tình huống 1.}
Một đỉnh có ít nhất ba khối con vừa liên quan vừa hoạt động ngoài. Do chỉ duyệt
hai hướng trên mặt ngoài, nhiều nhất hai khối như vậy có thể được xử lý trước
khi dừng. Khi xuất hiện ba khối, ta tìm được một \(\Kthree\) đồng phôi: ba
đỉnh phía trên là \((u,v,r)\), ba nhánh phía dưới đến từ ba khối con
\((a,b,c)\).

\begin{center}
  \small Hình 10: ba khối con hoạt động ngoài tạo thành một \(\Kthree\) đồng
  phôi.
\end{center}

\paragraph{Tình huống 2.}
Khi một đỉnh hoạt động ngoài \(w\) vẫn là đỉnh liên quan, thuật toán có thể đi
qua \(w\) xuống khối con. Việc này không làm mất \(w\) khỏi mặt ngoài: mọi cạnh
được nhúng ở một phía của khối con, và sau khi hợp nhất \(w\) không còn giữ
khối con đó nữa.

\begin{center}
  \small Hình 11: đi qua một đỉnh hoạt động ngoài để xử lý khối con của nó
  không che mất đỉnh đó khỏi mặt ngoài.
\end{center}

Nếu một khối con của \(w\) cần được duyệt hai lần thì đồ thị không phẳng. Có
ba trường hợp con:
\begin{enumerate}
  \item Trong khối con hoặc hậu duệ của nó có một đỉnh dừng tách hai đỉnh liên
  quan trở lên. Khi đó tồn tại một \(\Kthree\) đồng phôi, với hai phía tương
  ứng như \((u,a,c)\) và \((v,w,b)\).

  \item Trong khối con hoặc hậu duệ của \(w\) có hai đỉnh vừa liên quan vừa
  hoạt động ngoài. Khi đó tìm được một \(\Kfive\) đồng phôi.

  \item \(w\) có hai khối con vừa liên quan vừa hoạt động ngoài. Nếu mỗi khối
  chỉ cần một lần duyệt thì \(w\) vẫn nằm trên mặt ngoài; nếu có ba khối thì
  quay về tình huống 1; nếu một khối cần hai lần duyệt thì quay về hai trường
  hợp con trên.
\end{enumerate}

\begin{center}
  \small Hình 12 và 13: các cấu hình buộc một khối con phải duyệt hai lần lần
  lượt sinh ra \(\Kthree\) hoặc \(\Kfive\) đồng phôi.
\end{center}

\paragraph{Tình huống 3.}
Một đỉnh liên quan không thể được duyệt tới vì bị hai đỉnh dừng chắn lại.
Khi đó tồn tại một \(\Kthree\) đồng phôi, chẳng hạn bởi hai bộ nhánh
\((u,r,c)\) và \((v,a,b)\).

\begin{center}
  \small Hình 14: một đỉnh liên quan bị hai đỉnh dừng chặn dẫn tới
  \(\Kthree\) đồng phôi.
\end{center}

Như vậy mọi nguyên nhân làm quá trình nhúng thất bại đều chứng minh được đồ
thị không phẳng; còn nếu không thất bại, các cạnh được thêm trên mặt ngoài và
không cắt nhau. Thuật toán đúng.

\section{Bài liên quan và tổng kết}

Một số bài liên quan:
\begin{enumerate}
  \item Trại đông 2004: \textit{Bắp ngô tổ ong}.
  \item MIPT 090: \textit{Planarity graph}.
  \item Bài kiểm thử đi kèm: \textit{Bảng mạch}, kèm dữ liệu và mã nguồn.
\end{enumerate}

Thuật toán này có thể không xuất hiện thường xuyên trong thi đấu, nhưng việc
học nó giúp hiểu sâu hơn về đồ thị phẳng. Đối mặt với những bài khó hơn cũng
là cách rèn năng lực phân tích và cài đặt.

\section{Phụ lục}

Tài liệu tham khảo chính:
\begin{quote}
John M. Boyer, \textit{Simplified O(n) Planarity by Edge Addition}.
\end{quote}

Trong giả mã, ký hiệu \texttt{<-} là phép gán. Với đỉnh \(k\), \(k_{\rm in}\)
là hướng đi vào, \(k_{\rm out}\) là hướng đi ra. Trên mặt ngoài, hướng được mã
hóa bằng \(0\) và \(1\). Nếu một đỉnh có điểm sao chép, và khối của điểm sao
chép ấy có con duy nhất \(c\), ta lưu điểm sao chép dưới dạng \(c+n\).

\paragraph{Các hàm phụ.}
\begin{description}
  \item[\texttt{GetSuccessorOnExternalFace(v, vin)}] trả về đỉnh kế tiếp trên
  mặt ngoài.
  \item[\texttt{GetActiveSuccessorOnExternalFace(v, vin)}] trả về đỉnh hoạt
  động kế tiếp trên mặt ngoài.
  \item[\texttt{MergeBiconnectedComponent(S)}] hợp nhất các khối được lưu trên
  ngăn xếp, gồm cả xử lý phép đảo.
  \item[\texttt{EmbedBackEdge(v, vin, w, win)}] nhúng một cạnh hồi.
  \item[\texttt{Walkup(w, v)}] đi từ \(w\) lên \(v\).
  \item[\texttt{Walkdown(v)}] đi từ \(v\) xuống để nhúng các cạnh hồi.
\end{description}

\subsection*{\texttt{MergeBiconnectedComponent}}

\begin{verbatim}
Lay tu stack S bon gia tri: r, r_in, w, w_out
if r_in == w_out:
    hoan doi hai huong cua w
    dat canh goc (r, w) thanh -1
    w_out <- w_out xor 1

Xoa w khoi proots cua r
Xoa (w - n) khoi SDlist cua cha no

x <- 1 xor w_out
Dat huong v_in cua v tro toi x
Dat huong cua x dang tro toi w thanh tro toi v
\end{verbatim}

\subsection*{\texttt{walkup}}

Đang xử lý đỉnh \(V\), cần đi từ \(w\) lên \(V\).

\begin{verbatim}
Danh dau co canh hoi cua w bang V
(x, x_in) <- (w, 1)
(y, y_in) <- (w, 0)
while x != V:
    if visited[x] == V or visited[y] == V:
        break
    visited[x] <- V
    visited[y] <- V
    if x la goc:
        z1 <- x
    else if y la goc:
        z1 <- y
    else:
        z1 <- nil
    if z1 != nil:
        c <- z1 - n
        if lowpoint[c] < depth[V]:
            them z1 vao cuoi proots[z]
        else:
            them z1 vao dau proots[z]
    else:
        (x, x_in) <- GetSuccessorOnExternalFace(x, x_in)
        (y, y_in) <- GetSuccessorOnExternalFace(y, y_in)
\end{verbatim}

\subsection*{\texttt{walkdown}}

Đang xử lý đỉnh \(v\), đỉnh gốc được ký hiệu \(v_0\); cần đi xuống từ \(v_0\).

\begin{verbatim}
Xoa rong stack S
for v_out in {0, 1}:
    (w, w_in) <- GetSuccessorOnExternalFace(v)
    while w != v0:
        if backedge_flag[w] == v:
            while S khong rong:
                MergeBiconnectedComponent(S)
            EmbedBackEdge(v0, v_out, w, w_in)
            xoa backedge_flag[w]

        if proots[w] khong rong:
            day (w, w_in) vao S
            w0 <- phan tu dau cua proots[w]
            (x, x_in) <- GetActiveSuccessorOnExternalFace(w0, 0)
            (y, y_in) <- GetActiveSuccessorOnExternalFace(w0, 1)

            if x la dinh hoat dong trong:
                (w, w_in) <- (x, x_in)
            else if y la dinh hoat dong trong:
                (w, w_in) <- (y, y_in)
            else if x la dinh lien quan:
                (w, w_in) <- (x, x_in)
            else:
                (w, w_in) <- (y, y_in)

            if w == x:
                w_out <- 0
            else:
                w_out <- 1
            day (w0, w_out) vao S

        else if w la dinh hoat dong trong:
            (w, w_in) <- GetSuccessorOnExternalFace(w, w_in)
        else:
            if lowpoint[c] < depth[v] and S rong:
                nhung canh gia (v0, w)
            break

    if S khong rong:
        break
\end{verbatim}

\end{document}
